% Document/LinearAlgebra/VectorSpaces/Matrices/MatrixSpaces.tex
% Phase 16a–d: Matrix Spaces, Representation Theorem, Operations by Transport, Action Map

\section{Matrix Spaces}
\label{sec:mat-spaces}

% ── Motivation ──────────────────────────────────────────────────────────────

\begin{remark}[Why Matrices?]
    We have built a rich theory of linear maps. But working with abstract maps is cumbersome
    for computation: to describe $T \in \Hom[K]{U}{V}$ we would need to specify $T(v)$ for
    every $v \in U$. By the universal property, $T$ is completely determined by its values
    on a basis $\mathcal{B}_U = (b_j)_{j \in J}$, reducing the data to the family
    $(T(b_j))_{j \in J}$ in $V$. Expressing each $T(b_j)$ in a basis
    $\mathcal{B}_V = (c_i)_{i \in I}$ of $V$ further reduces the information to a family
    of scalars $(a_{ij})_{i \in I, j \in J}$ with $T(b_j) = \sum_i a_{ij} c_i$.
    A \textbf{matrix} is precisely this rectangular array of scalars. It stores the action
    of $T$ in a form that is:
    \begin{itemize}
        \item \textbf{Finite}: finitely many nonzero entries per column, since $T(b_j) \in V$
              has finite support in $\mathcal{B}_V$,
        \item \textbf{Computable}: finding $T(v)$ for any $v$ reduces to arithmetic on scalars,
        \item \textbf{Composable}: the matrix of $\Composition{S}{T}$ is determined by the
              matrices of $S$ and $T$.
    \end{itemize}
    The array depends on the choice of ordered bases. We begin by formalizing the existence of
    this array, then define matrices as purely combinatorial objects (grids of scalars). The
    \textbf{Representation Theorem} (Theorem~\ref{thm:rep-thm}) establishes that the
    column-finite ones are \emph{exactly} the linear maps between free modules with canonical
    bases. Matrix operations are then defined by \textbf{transport of structure} from
    composition of linear maps, and the classical \textbf{commutative diagram}
    (Section~\ref{sec:matrix-of-map}) characterizes the matrix of an abstract linear map.
\end{remark}

\begin{proposition}[Existence and Uniqueness of the Scalar Array]
\label{prop:scalar-array}
    Let $T: U \to V$ be $K$-linear, with ordered bases $\mathcal{B}_U = (b_j)_{j \in J}$
    and $\mathcal{B}_V = (c_i)_{i \in I}$. For each $j \in J$ there exist \textbf{unique}
    scalars $(A_{ij})_{i \in I}$ with $A_{ij} \in K$ such that
    \[
        T(b_j) = \sum_{i \in I} A_{ij}\, c_i
    \]
    and for each fixed $j$, only finitely many $A_{ij}$ are nonzero.
\end{proposition}

\begin{proof}
    Each $T(b_j) \in V$. Since $\mathcal{B}_V$ is a basis of $V$, every vector in $V$ has a
    unique expression as a finite linear combination of $(c_i)_{i \in I}$. Applying this to
    $T(b_j)$ gives the unique scalars $A_{ij} = \CoordVec[\mathcal{B}_V]{T(b_j)}(i)$, with
    only finitely many nonzero (finite support by definition of a basis).
\end{proof}

\begin{remark}[The Fundamental Datum]
    The double index $(i,j)$ is the fundamental datum of this chapter: $i$ picks the row
    (the target basis vector $c_i$) and $j$ picks the column (the source basis vector $b_j$).
    The family $(A_{ij})_{i \in I,\, j \in J}$ is the \textbf{matrix of $T$} with respect to
    $\mathcal{B}_U$ and $\mathcal{B}_V$. Everything that follows --- evaluation formula,
    matrix-times-vector, matrix multiplication --- is a consequence of this single fact
    together with linearity. The key question is: \emph{what is a matrix, as a mathematical
    object, independent of the linear map that produced it?}
\end{remark}

% ── 16a: Definitions ────────────────────────────────────────────────────────

\subsection{General Matrices}

\begin{definition}[{General Matrix Space $\MatSpace{Y}{X}[K]$}]
\label{def:mat-space}
    Let $X, Y$ be sets (column and row index sets, respectively). A \textbf{$(Y \times X)$-matrix
    with entries in $K$} is a function
    \[
        A: Y \times X \to K
    \]
    Write $A_{ij}$ for $A(i,j)$ (row index $i \in Y$, column index $j \in X$). The set of
    all such matrices is
    \[
        \MatSpace{Y}{X}[K] \defeq K^{Y \times X}
    \]
    with no finiteness condition imposed. For cardinals $\kappa_1, \kappa_2$, write
    $\MatSpace{\kappa_2}{\kappa_1}[K]$ ($\kappa_2$ rows, $\kappa_1$ columns).
\end{definition}

\begin{remark}[Scope]
    The general space $\MatSpace{Y}{X}[K]$ imposes no finiteness condition. In this chapter,
    most results concern the column-finite subspace $\CFMat{Y}{X}[K]$, which captures exactly
    the linear maps $K^{(\kappa_1)} \to K^{(\kappa_2)}$
    (Theorem~\ref{thm:rep-thm}). The larger space $\MatSpace{Y}{X}[K]$ becomes essential in
    functional analysis, where bounded operators on sequence spaces are represented by infinite
    matrices with no column-finiteness requirement --- convergence replaces finite support as
    the operative condition.
\end{remark}

\begin{remark}[Currying]
    A matrix $A \in \MatSpace{Y}{X}[K]$ is equivalently a function
    $j \mapsto A(-,j) \in K^Y$ (the column map) or $i \mapsto A(i,-) \in K^X$ (the row map).
\end{remark}

\begin{definition}[{Matrix Addition and Scalar Multiplication on $\MatSpace{Y}{X}[K]$}]
\label{def:mat-add-scalar}
    Since $\MatSpace{Y}{X}[K] = K^{Y \times X}$ is a product of copies of $K$, it inherits
    pointwise operations making it a $K$-vector space:
    \[
        (A + B)_{ij} \defeq A_{ij} + B_{ij}, \qquad
        (\lambda A)_{ij} \defeq \lambda\, A_{ij}, \qquad
        0_{ij} \defeq 0
    \]
    These are defined for \textbf{all} $A, B \in \MatSpace{Y}{X}[K]$, with no finiteness
    condition required.
\end{definition}

\subsection{Column-Finite Matrices}

\begin{definition}[{Column-Finite Matrix Space $\CFMat{Y}{X}[K]$}]
\label{def:cfmat-space}
    A matrix $A \in \MatSpace{Y}{X}[K]$ is \textbf{column-finite} if for each $j \in X$,
    the column $A(-,j): i \mapsto A_{ij}$ has \textbf{finite support}:
    $A_{ij} = 0$ for all but finitely many $i \in Y$. Equivalently, $A(-,j) \in K^{(Y)}$
    for all $j$. The set of column-finite matrices is
    \[
        \CFMat{Y}{X}[K] \defeq \Set{A \in \MatSpace{Y}{X}[K] \mid
            \forall j \in X,\ A(-,j) \in K^{(Y)}}
    \]
\end{definition}

\begin{remark}[Why Column-Finiteness?]
    Column-finiteness is not an arbitrary restriction. When expressing a linear map $T: U \to V$
    in ordered bases, each image $T(b_j) \in V$ is a \emph{finite} linear combination of the
    target basis vectors --- this is a consequence of the definition of a basis. Column-finiteness
    of the matrix is precisely this condition on the scalars.
\end{remark}

\begin{proposition}[{$\CFMat{Y}{X}[K]$ is a Subspace}]
\label{prop:cfmat-subspace}
    $\CFMat{Y}{X}[K]$ is a subspace of $\MatSpace{Y}{X}[K]$, and therefore a $K$-vector space
    under the operations of Definition~\ref{def:mat-add-scalar}.
\end{proposition}

\begin{proof}
    The zero matrix has $0(-,j) = 0 \in K^{(Y)}$ for all $j$.
    For $A, B \in \CFMat{Y}{X}[K]$ and $\lambda \in K$:
    $(A + B)(-,j) = A(-,j) + B(-,j) \in K^{(Y)}$ (sum of finitely supported functions
    is finitely supported) and $(\lambda A)(-,j) = \lambda\, A(-,j) \in K^{(Y)}$.
\end{proof}

\begin{remark}[Currying for Column-Finite Matrices]
    Via currying, $\CFMat{Y}{X}[K] \cong (K^{(Y)})^X$ --- functions from $X$ to
    finitely-supported column vectors. This is \textbf{not} the same as $K^{(Y \times X)}$:
    the latter requires globally finite support (only finitely many nonzero entries in total),
    which is strictly stronger than column-finiteness when $X$ is infinite. A column-finite
    matrix may have infinitely many nonzero columns, each individually finitely supported.
\end{remark}

\subsection{Sizes}

\begin{remark}[General vs.\ Column-Finite]
    \leavevmode
    \begin{itemize}
        \item \textbf{General matrices}: $\MatSpace{Y}{X}[K] = K^{Y \times X}$, canonically
              isomorphic to $\prod_{j \in X} K^Y$ (one copy of $K^Y$ per column).
        \item \textbf{Column-finite matrices}: $\CFMat{Y}{X}[K] \cong (K^{(Y)})^X$ --- functions
              from $X$ to finitely-supported column vectors.
        \item \textbf{Contrast with $K^{(Y \times X)}$}: The space $K^{(Y \times X)}$ requires
              globally finite support (only finitely many nonzero entries in total). This is
              strictly stronger than column-finiteness when $X$ is infinite.
    \end{itemize}
\end{remark}

% ── 16b: The Representation Theorem ────────────────────────────────────────

\newpage
\subsection{The Representation Theorem}

Write $\CanonicalBase{\kappa} = (\CanonicalVector{i})_{i \in \kappa}$ for the canonical ordered basis of
$K^{(\kappa)}$.

\begin{theorem}[{Representation Theorem: $\Hom[K]{K^{(\kappa_1)}}{K^{(\kappa_2)}} \cong \CFMat{\kappa_2}{\kappa_1}[K]$}]
\label{thm:rep-thm}
    There is a canonical $K$-linear isomorphism
    \[
        \Phi: \Hom[K]{K^{(\kappa_1)}}{K^{(\kappa_2)}} \xrightarrow{\;\sim\;}
        \CFMat{\kappa_2}{\kappa_1}[K]
    \]
    defined by
    \[
        \Phi(T)_{ij} \defeq \bigl(T(\CanonicalVector{j})\bigr)_i
    \]
    --- the $i$-th coordinate of $T(\CanonicalVector{j})$ in $K^{(\kappa_2)}$.
    Equivalently, the $j$-th column of $\Phi(T)$ is $T(\CanonicalVector{j}) \in K^{(\kappa_2)}$.
    The codomain is $\CFMat{\kappa_2}{\kappa_1}[K]$ (not the larger
    $\MatSpace{\kappa_2}{\kappa_1}[K]$): column-finiteness is not imposed externally but is a
    \emph{consequence} of $T$ mapping into $K^{(\kappa_2)}$.
\end{theorem}

\begin{remark}[Convention: Why $(T(\CanonicalVector{j}))_i$ and Not $(T(\CanonicalVector{i}))_j$?]
    The definition $\Phi(T)_{ij} = (T(\CanonicalVector{j}))_i$ is a choice: we could equally
    define the $(i,j)$-entry as $(T(\CanonicalVector{i}))_j$, which would identify
    $\Hom[K]{K^{(\kappa_1)}}{K^{(\kappa_2)}}$ with $\CFMat{\kappa_1}{\kappa_2}[K]$ (a
    $\kappa_1 \times \kappa_2$ matrix) rather than $\CFMat{\kappa_2}{\kappa_1}[K]$ (a
    $\kappa_2 \times \kappa_1$ matrix). In some sense this alternative is more natural: a map
    $K^{(\kappa_1)} \to K^{(\kappa_2)}$ would produce a matrix whose dimensions match the
    order $(\kappa_1, \kappa_2)$.

    The classical convention $(T(\CanonicalVector{j}))_i$ reverses these dimensions, but has the compensating
    advantage that matrix multiplication follows the familiar left-to-right reading order:
    a $\kappa_2 \times \kappa_1$ matrix times a $\kappa_1 \times \kappa_3$ matrix yields a
    $\kappa_2 \times \kappa_3$ matrix, with adjacent indices contracting. Under the alternative
    convention, the same contraction would require the reverse order. We follow the classical
    convention throughout.
\end{remark}

\begin{proof}
    We construct a $K$-linear inverse $\Psi: \CFMat{\kappa_2}{\kappa_1}[K] \to
    \Hom[K]{K^{(\kappa_1)}}{K^{(\kappa_2)}}$.

    \textit{$\Phi$ is well-defined and $K$-linear}:
    Each $T(\CanonicalVector{j}) \in K^{(\kappa_2)}$ has finite support in $i$, so
    $\Phi(T)(-,j) \in K^{(\kappa_2)}$, i.e.\ the matrix is column-finite. Linearity:
    $\Phi(\alpha T + \beta S)_{ij}
    = ((\alpha T + \beta S)(\CanonicalVector{j}))_i
    = \alpha\, \Phi(T)_{ij} + \beta\, \Phi(S)_{ij}$.

    \textit{Construction of $\Psi$}:
    Given $A \in \CFMat{\kappa_2}{\kappa_1}[K]$, define $\Psi(A) = T_A$ by
    $T_A(\CanonicalVector{j}) \defeq A(-,j) \in K^{(\kappa_2)}$ (well-defined since $A$ is
    column-finite) and extend by linearity. The map $A \mapsto T_A$ is $K$-linear
    (since $(A + B)(-,j) = A(-,j) + B(-,j)$ and $(\lambda A)(-,j) = \lambda\, A(-,j)$).

    \textit{$\Psi$ is a two-sided inverse}:
    $(\Phi \circ \Psi)(A)_{ij} = (T_A(\CanonicalVector{j}))_i = A(-,j)(i) = A_{ij}$,
    so $\Phi \circ \Psi = \Identity{\CFMat{\kappa_2}{\kappa_1}[K]}$.
    $(\Psi \circ \Phi)(T)(\CanonicalVector{j}) = \Phi(T)(-,j) = T(\CanonicalVector{j})$,
    so $\Psi(\Phi(T))$ and $T$ agree on the basis $\CanonicalBase{\kappa_1}$, hence
    $\Psi \circ \Phi = \Identity{\Hom[K]{K^{(\kappa_1)}}{K^{(\kappa_2)}}}$.

    \textit{Matrix-times-column formula}:
    For any $x = \sum_j x_j \CanonicalVector{j} \in K^{(\kappa_1)}$ (a finite sum):
    \[
        T_A(x) = \sum_j x_j\, A(-,j), \qquad
        (T_A(x))_i = \sum_j A_{ij}\, x_j = (Ax)_i
    \]
    so $T_A(x) = Ax$. Every $A \in \CFMat{\kappa_2}{\kappa_1}[K]$ acts on
    $K^{(\kappa_1)} \to K^{(\kappa_2)}$ by left multiplication.
\end{proof}

\begin{definition}[{Identity Matrix $\IdentityMat{\kappa}$}]
\label{def:identity-matrix}
    The \textbf{identity matrix} $\IdentityMat{\kappa} \in \CFMat{\kappa}{\kappa}[K]$ is
    the matrix with entries
    \[
        (\IdentityMat{\kappa})_{ij} \defeq \delta_{ij}
        = \begin{cases} 1 & \text{if } i = j \\ 0 & \text{if } i \neq j \end{cases}
    \]
    For finite $n$, write $\IdentityMat{n}$ for the $n \times n$ identity matrix.
\end{definition}

\begin{corollary}[Special Matrices]
\label{cor:special-matrices}
    Under the isomorphism $\Phi$:
    \begin{enumerate}
        \item \textbf{Zero matrix = zero map}:
              $\Phi(0)_{ij} = (0 \cdot \CanonicalVector{j})_i = 0$.
        \item \textbf{Identity matrix = identity map}:
              $\Phi(\Identity{K^{(\kappa)}}) = \IdentityMat{\kappa}$.

              \textit{Proof}: $\Phi(\Identity{K^{(\kappa)}})_{ij}
              = (\Identity{K^{(\kappa)}}(\CanonicalVector{j}))_i = (\CanonicalVector{j})_i = \delta_{ij}
              = (\IdentityMat{\kappa})_{ij}$.
        \item \textbf{Scalar matrix}:
              $\Phi(\lambda\, \Identity{K^{(\kappa)}})
              = \lambda\, \IdentityMat{\kappa}$.
    \end{enumerate}
\end{corollary}

\begin{remark}[{$V \cong \Hom[K]{K}{V}$ as a Special Case}]
\label{rem:vec-as-hom}
    The isomorphism $\Hom[K]{K}{V} \cong V$ of Theorem~\ref{thm:hom-k-v-iso} is the
    $\kappa_1 = 1$ special case of the Representation Theorem:
    $\CFMat{\kappa_2}{1}[K] \cong \Hom[K]{K}{K^{(\kappa_2)}} \cong K^{(\kappa_2)}$.
    Column matrices with $\kappa_2$ rows \textbf{are} vectors in $K^{(\kappa_2)}$.
\end{remark}

\begin{remark}[{Grid vs.\ Map: Two Structures on the Same Set}]
\label{rem:grid-vs-map}
    The Representation Theorem establishes a canonical bijection
    $\Phi: \Hom[K]{K^{(\kappa_1)}}{K^{(\kappa_2)}} \xrightarrow{\sim}
    \CFMat{\kappa_2}{\kappa_1}[K]$. The domain and codomain are the \emph{same underlying
    set} --- a column-finite matrix $A$ and the linear map $x \mapsto Ax$ carry identical
    information. But they support different operations:
    \begin{itemize}
        \item \textbf{As an element of
              $\Hom[K]{K^{(\kappa_1)}}{K^{(\kappa_2)}}$} (the ``map'' view): $A$ can be
              composed with other maps, it has a kernel and an image, and it transforms under
              change of basis by conjugation. These are the \emph{$K$-linear} operations ---
              they see $A$ as a morphism in the category $K$-Vect.
        \item \textbf{As an element of
              $\CFMat{\kappa_2}{\kappa_1}[K]$} (the ``grid'' view): $A$ is a function
              $\kappa_2 \times \kappa_1 \to K$, and we can perform entrywise operations ---
              the Hadamard product $(\Hadamard{A}{B})_{ij} = A_{ij} B_{ij}$, entrywise
              absolute value, transpose as the index swap $(i,j) \mapsto (j,i)$, extraction of
              rows/columns/diagonals, sparsity analysis, reshaping. None of these are natural
              operations on $\operatorname{Hom}$.
    \end{itemize}
    This is the same phenomenon as for vectors
    (Remark~\ref{rem:three-layers-vector}): $K^{(I)}$ is simultaneously a vector space and a
    set of $I$-tuples. The coordinate isomorphism bridges the geometric world to the
    computational world; the Representation Theorem does the same for linear maps. We identify
    the two views via $\Phi$ and use context to determine which structure is active.

    \textbf{Convention}: From this point on, we use $K$-linear maps
    $K^{(\kappa_1)} \to K^{(\kappa_2)}$ and column-finite matrices in
    $\CFMat{\kappa_2}{\kappa_1}[K]$ interchangeably via $\Phi$, without further comment.
    In particular, we write $A$ both for a matrix and for the linear map $\Phi^{-1}(A)$,
    and $T$ both for a linear map and for the matrix $\Phi(T)$.

    For a linear map $T: U \to V$ between \emph{abstract} vector spaces (with no canonical
    basis), a dedicated notation $\CoordMatrix[\mathcal{B}_U][\mathcal{B}_V]{T}$ for its
    matrix with respect to ordered bases $\mathcal{B}_U$, $\mathcal{B}_V$ will be
    introduced in Section~\ref{sec:matrix-of-map}.

    This is the same pattern at three levels of abstraction:

    \medskip
    \begin{center}
    \begin{tabular}{l|lll}
    \hline
    \textbf{Level} & \textbf{Geometric object} & \textbf{Canonical representation} & \textbf{Data structure} \\
    \hline
    Scalar & $a \in K$ & $a \in K$ & a field element \\[4pt]
    Vector & $v \in V$ & $\Chart{\mathcal{B}}(v) \in K^{(I)}$ & tuple $(a_i)_{i \in I}$ \\[4pt]
    Map & $T \in \Hom[K]{U}{V}$ &
        $\Chart{\mathcal{B}_V} \circ T \circ \Frame{\mathcal{B}_U}$ &
        grid $(A_{ij})$ \\
        & & $\in \Hom[K]{K^{(\kappa_1)}}{K^{(\kappa_2)}}$ &
        $(i,j) \in \kappa_2 \times \kappa_1$ \\
    \hline
    \end{tabular}
    \end{center}
    \medskip

    At each level, the canonical representation is a vector space element (in a space with a
    canonical basis), and the data structure is the same set viewed as a family of scalars.
    The representation theorem ($\Phi$) is the bridge between the last two columns at the Map
    level, just as coordinate extraction is the bridge at the Vector level.
\end{remark}

% ── The Action Map ─────────────────────────────────────────────────────────

\subsection{The Action Map}

\begin{proposition}[Action Map]
\label{prop:action-map}
    Every $A \in \MatSpace{Y}{X}[K]$ defines a $K$-linear map
    \[
        \mu_A: K^{(X)} \to K^Y, \qquad
        (\mu_A(x))_i \defeq \sum_{j \in X} A_{ij}\, x_j
    \]
    The sum is finite (since $x$ has finite support), so $\mu_A$ is well-defined.
    Column-finiteness of $A$ is equivalent to
    $\operatorname{im}(\mu_A) \subseteq K^{(Y)}$: i.e.\ $A$ is column-finite if and only if
    $\mu_A$ restricts to a map $K^{(X)} \to K^{(Y)}$.
\end{proposition}

\begin{proof}
    If $A$ is column-finite, then for any $x \in K^{(X)}$,
    $\mu_A(x) = \sum_{j \in \operatorname{supp}(x)} x_j\, A(-,j)$ is a finite sum of
    elements of $K^{(Y)}$, so $\mu_A(x) \in K^{(Y)}$. Conversely, if
    $\mu_A(\CanonicalVector{j}) = A(-,j) \in K^{(Y)}$ for all $j$, then $A$ is column-finite.
\end{proof}

\begin{remark}[Scope]
    The action map $\mu_A$ is well-defined for \textbf{all} matrices in $\MatSpace{Y}{X}[K]$,
    not only column-finite ones --- the sum $\sum_j A_{ij} x_j$ is finite because $x$ has
    finite support. What column-finiteness buys is that the result lands in $K^{(Y)}$ rather
    than merely in $K^Y$. In functional analysis, convergence replaces finite support as the
    operative condition, but this is not our concern here: we work primarily with
    finite-dimensional spaces (for which $\CFMat{m}{n}[K] = \MatSpace{m}{n}[K]$ and all
    conditions are automatic).

    One cannot extend the domain of matrix-vector multiplication from $K^{(X)}$ to $K^X$
    without a topology: if $x \in K^X$ has infinite support, the sum $\sum_j A_{ij} x_j$ may
    have infinitely many nonzero terms, and there is no purely algebraic way to assign it a
    value. The passage from finite support (algebraic) to convergent sums (analytic) is exactly
    where functional analysis begins. In that setting, coordinate representations use
    topological bases (Schauder bases, orthonormal bases) rather than Hamel bases, the
    coordinate map targets a space like $\ell^2(I)$ rather than $K^{(I)}$, and the resulting
    ``matrices'' need not be column-finite --- they must instead satisfy boundedness conditions
    dictated by the topology. A different notation for these topological coordinate
    representations is warranted, since the target space and the governing finiteness conditions
    are genuinely different.
\end{remark}

% ── 16c: Matrix Operations by Transport of Structure ───────────────────────

\newpage
\subsection{Matrix Multiplication}

\begin{remark}[Transport of Structure]
    Matrix addition and scalar multiplication on $\CFMat{\kappa_2}{\kappa_1}[K]$ are
    compatible with $\Phi$: the $K$-linearity of $\Phi$ (proved in
    Theorem~\ref{thm:rep-thm}) means that the pointwise formulas of
    Definition~\ref{def:mat-add-scalar} are the \emph{unique} operations making $\Phi$ a
    $K$-linear map.

    Matrix multiplication has no natural a priori definition; it is \emph{forced upon us} by
    requiring $\Phi$ to be an algebra morphism. Since $\End[K]{K^{(\kappa)}}$ is a $K$-algebra
    with composition as multiplication, requiring
    $\Phi(\Composition{S}{T}) = \Phi(S) \cdot \Phi(T)$ uniquely determines the product of
    column-finite matrices.
\end{remark}

\begin{definition}[Matrix Multiplication]
\label{def:mat-mult}
    Let $A \in \CFMat{Y}{Z}[K]$ and $B \in \CFMat{Z}{X}[K]$. Define
    \[
        AB \defeq \Phi\bigl(\Composition{\Phi^{-1}(A)}{\Phi^{-1}(B)}\bigr)
    \]
    That is, $AB$ is the unique matrix corresponding to the composition of the associated
    linear maps.
\end{definition}

\begin{theorem}[Entry Formula]
\label{thm:entry-formula}
    For $A \in \CFMat{Y}{Z}[K]$ and $B \in \CFMat{Z}{X}[K]$:
    \[
        (AB)_{ij} = \sum_{k \in Z} A_{ik}\, B_{kj}
    \]
\end{theorem}

\begin{proof}
    Let $T_A = \Phi^{-1}(A)$ and $T_B = \Phi^{-1}(B)$. The $j$-th column of $AB$ is:
    \[
        \Phi(\Composition{T_A}{T_B})(-,j)
        = (\Composition{T_A}{T_B})(\CanonicalVector{j})
        = T_A(T_B(\CanonicalVector{j}))
        = T_A(B(-,j))
    \]
    Since $B(-,j) = \sum_k B_{kj}\, \CanonicalVector{k} \in K^{(Z)}$ (finite sum by column-finiteness of
    $B$), linearity of $T_A$ gives:
    \[
        T_A(B(-,j)) = \sum_k B_{kj}\, T_A(\CanonicalVector{k}) = \sum_k B_{kj}\, A(-,k)
    \]
    The $i$-th entry is $(AB)_{ij} = \sum_{k \in Z} A_{ik}\, B_{kj}$.
\end{proof}

\begin{proposition}[Well-Definedness of Matrix Multiplication]
\label{prop:mat-mult-wd}
    The product $AB$ lies in $\CFMat{Y}{X}[K]$.
\end{proposition}

\begin{proof}
    Fix $j \in X$. Let $S_j = \{k \in Z : B_{kj} \neq 0\}$; this set is finite by
    column-finiteness of $B$. Then
    $(AB)_{ij} = \sum_{k \in S_j} A_{ik}\, B_{kj}$
    is a finite sum in $K$, well-defined for all $i \in Y$.
    Moreover, $\{i \in Y : (AB)_{ij} \neq 0\}
    \subseteq \bigcup_{k \in S_j} \{i \in Y : A_{ik} \neq 0\}$,
    which is a finite union of finite sets (by column-finiteness of $A$),
    so $(AB)(-,j) \in K^{(Y)}$.
\end{proof}

\begin{remark}[When Is Multiplication Defined for General Matrices?]
    The entry formula $(AB)_{ij} = \sum_k A_{ik} B_{kj}$ requires the sum to be finite.
    This is guaranteed when $B$ is column-finite (then $B_{kj} \neq 0$ for finitely many $k$).
    So one can multiply $A \in \MatSpace{Y}{Z}[K]$ by $B \in \CFMat{Z}{X}[K]$ to get
    $AB \in \MatSpace{Y}{X}[K]$ --- but the result need not be column-finite (this requires
    $A$ to also be column-finite). Column-finiteness of both factors is the minimal condition
    ensuring the product is well-defined \emph{and} stays in $\CFMat{Y}{X}[K]$.
\end{remark}

\begin{remark}[Properties Inherited from Composition]
    Since $AB$ is defined as $\Phi(\Composition{\Phi^{-1}(A)}{\Phi^{-1}(B)})$ and $\Phi$
    is a $K$-linear isomorphism, all algebraic properties of matrix multiplication are
    inherited from composition of linear maps \textbf{without any index calculations}:
    \begin{itemize}
        \item \textbf{Associativity}: $(AB)C = A(BC)$, from associativity of composition.
        \item \textbf{Bilinearity}: $A(B + C) = AB + AC$,\;
              $(A + B)C = AC + BC$,\;
              $\lambda(AB) = (\lambda A)B = A(\lambda B)$, from linearity.
        \item \textbf{Identity}: $\IdentityMat{\kappa} \cdot A = A$ and $A \cdot \IdentityMat{\kappa} = A$,
              from $\Composition{\Identity{}}{T} = T = \Composition{T}{\Identity{}}$.
        \item \textbf{Zero}: $0 \cdot A = 0$ and $A \cdot 0 = 0$, from composing with the
              zero map.
    \end{itemize}
\end{remark}

\begin{corollary}[{$K$-Algebra Isomorphism $\End[K]{K^{(\kappa)}} \cong \CFMat{\kappa}{\kappa}[K]$}]
\label{cor:end-algebra-iso}
    $\Phi$ is an isomorphism of $K$-algebras:
    \begin{itemize}
        \item Adding maps $\leftrightarrow$ adding matrices
        \item Composing maps $\leftrightarrow$ multiplying matrices
        \item Zero map $\leftrightarrow$ zero matrix
        \item Identity map $\leftrightarrow$ identity matrix
    \end{itemize}
\end{corollary}

\subsection{Other Matrix Operations --- Grid-Side Operations}

\begin{remark}[Grid-Side Operations]
    The operations in this subsection are \emph{grid-side} operations: they act on entries of the
    matrix viewed as a function $Y \times X \to K$. They do not, in general, correspond to natural
    operations on $\Hom[K]{K^{(X)}}{K^{(Y)}}$. The Hadamard product, for instance, depends on the
    choice of canonical basis and has no intrinsic meaning for abstract linear maps. This is the
    grid/map distinction from Remark~\ref{rem:grid-vs-map} in action.
\end{remark}

\begin{definition}[Hadamard Product]
\label{def:hadamard}
    For $A, B \in \MatSpace{Y}{X}[K]$, the \textbf{Hadamard product} (entrywise product) is
    \[
        (\Hadamard{A}{B})_{ij} \defeq A_{ij}\, B_{ij}
    \]
    This is defined on all of $\MatSpace{Y}{X}[K]$ (same-size matrices only).
    If $A, B \in \CFMat{Y}{X}[K]$, then $\Hadamard{A}{B} \in \CFMat{Y}{X}[K]$
    (column-finiteness is preserved: $\{i : A_{ij} B_{ij} \neq 0\}
    \subseteq \{i : A_{ij} \neq 0\}$).

    This is a grid operation: it multiplies corresponding \emph{entries}, not maps. It depends
    on the indexing and has no natural interpretation in $\operatorname{Hom}$.
\end{definition}

\begin{definition}[Transpose]
\label{def:transpose}
    For $A \in \MatSpace{Y}{X}[K]$, the \textbf{transpose} $\Transpose{A} \in \MatSpace{X}{Y}[K]$
    is defined by
    \[
        (\Transpose{A})_{ij} \defeq A_{ji}
    \]
\end{definition}

\begin{remark}[Transpose and Column-Finiteness]
    The transpose of a column-finite matrix $A \in \CFMat{Y}{X}[K]$ is \textbf{not} in general
    column-finite as an element of $\MatSpace{X}{Y}[K]$: transposing turns columns of $A$ into
    rows of $\Transpose{A}$, and each row of $\Transpose{A}$ may have infinite support even if
    each column of $A$ does not. For finite matrices (both $X$ and $Y$ finite),
    $\Transpose{A} \in \CFMat{X}{Y}[K]$ always.
\end{remark}

\begin{remark}[Transpose: Grid vs.\ Algebra]
    As a grid operation, transposing simply swaps indices. Its relationship to the
    \emph{algebraic} transpose (the dual map $T^*: V^* \to U^*$) is a theorem, not a
    tautology --- the algebraic transpose is defined basis-independently via
    $T^*(\varphi) = \Composition{\varphi}{T}$, and the fact that its matrix is the
    grid-transpose of the matrix of $T$ requires proof using dual bases.
\end{remark}

\begin{definition}[Hermitian Conjugate]
\label{def:hermitian}
    If $K$ carries a conjugation involution $\overline{(\cdot)}: K \to K$ (e.g.\ $K = \bb{C}$),
    the \textbf{Hermitian conjugate} (conjugate transpose) is
    \[
        (\Hermitian{A})_{ij} \defeq \overline{A_{ji}}
    \]
    Same column-finiteness caveat as for transpose.
\end{definition}

\begin{remark}[Determinant --- Forward Reference]
    For square finite matrices $A \in \CFMat{n}{n}[K]$ ($n < \infty$) over a commutative ring,
    the \textbf{Leibniz formula}
    \[
        \Det{A} = \sum_{\sigma \in S_n} \operatorname{sgn}(\sigma) \prod_{i=1}^{n} A_{i,\sigma(i)}
    \]
    will be derived in Chapter~7 (Multilinear Algebra) as the unique alternating multilinear form
    normalized by $\Det{\IdentityMat{n}} = 1$. For now, use $\Det{A}$ as notation and rely on the Leibniz
    formula when needed for small matrices.
\end{remark}
